\section{Conclusion}
We introduce \method, a plug-and-play multi-agent framework that 
coordinates the memory cycle along its forward and backward paths.
On the forward path, a Meta-Thinker separates strategic reasoning 
from low-level execution, addressing strategic blindness in 
construction and retrieval.
On the backward path, in-situ self-evolution converts probe QA 
failures into direct memory repair before the memory is committed.
Experiments on LoCoMo show that \method outperforms all baselines 
across multiple backbones and consistently improves three different 
storage backends.

\section{Limitations}
Our evaluation focuses on a dialogue-centric long-horizon memory benchmark.
While LoCoMo covers diverse question types, including single-hop, multi-hop, temporal, and open-domain reasoning, it does not capture all settings in which persistent memory may be needed.

% Second, probe generation currently relies on a strong LLM (Claude-Opus-4.5~\cite{anthropic2025claudeopus45}); while Appendix~\ref{appendix:impact_probe_generation_model} shows that even lighter models provide meaningful repair, the use of open-source models for probe generation remains unexplored.

In addition, the backward path assumes that interaction streams can be organized into sessions and that synthetic probe QA can provide useful localized supervision.
These assumptions are natural for the benchmark studied here, but may require adaptation in settings with less clear session boundaries or more open-ended interaction structure.

\section{Ethics Statement}
This work studies long-horizon memory management for LLM agents.
All experiments are conducted on the publicly available benchmark, which consists of synthetic conversations and does not contain real user data.
No personally identifiable information is collected, stored, or processed in this work.
We note that improving memory quality in agent systems may raise broader considerations for real-world deployment, including user privacy, informed consent for data retention, controllability over stored memories, and the risk of persisting incorrect information through automated repair.
While these concerns are beyond the scope of the present study, we believe they should be treated as first-class design requirements in any production deployment of memory-augmented agents.
